\section{System requierments}
\begin{tcolorbox}
  [colback=red!5!white,colframe=red!75!black,
  title=Note]
  This chapter outlines the system requirements for the proposed Type3 Asset Administration Shell (AAS) integrated with a Multi-Agent System (MAS) using the Belief-Desire-Intention (BDI) approach. It details both functional and non-functional requirements, providing a comprehensive overview of the necessary specifications to ensure successful implementation and operation of the system.
\end{tcolorbox}

\subsection{Functional requirements}
\begin{tcolorbox}[colback=red!5!white,colframe=red!75!black,
  title=About Functional Requirements]
  This section outlines the functional requirements for the Type3 Asset Administration Shell (AAS) integrated with a Multi-Agent System (MAS) using the Belief-Desire-Intention (BDI) approach. It specifies the essential functionalities that the system must possess to effectively manage and represent industrial assets, facilitate communication and cooperation among agents, and support autonomous decision-making processes. The requirements detailed in this section serve as a foundation for the design and implementation of the proposed system.
\end{tcolorbox}

\subsection{Non-functional requirements}
\begin{tcolorbox}[colback=red!5!white,colframe=red!75!black,
  title=About Non-Functional Requirements]
  This section delineates the non-functional requirements for the Type3 Asset Administration Shell (AAS) integrated with a Multi-Agent System (MAS) using the Belief-Desire-Intention (BDI) approach. It addresses critical aspects such as performance, scalability, security, and usability that are essential for the effective operation and adoption of the system in industrial environments. The non-functional requirements outlined in this section ensure that the system not only meets its functional objectives but also adheres to quality standards necessary for real-world applications.
\end{tcolorbox}
\subsection{Usecase}
\begin{tcolorbox}[colback=red!5!white,colframe=red!75!black,
  title=About Usecase]
  This section presents a detailed use case scenario for the Type3 Asset Administration Shell (AAS) integrated with a Multi-Agent System (MAS) using the Belief-Desire-Intention (BDI) approach. It illustrates how the system operates in a practical industrial context, highlighting the interactions between the AAS and BDI agents to achieve specific objectives. The use case serves to demonstrate the application of the functional and non-functional requirements outlined in previous sections, providing a concrete example of the system's capabilities and benefits.
\end{tcolorbox}